\section{Introduction: The Phenomenology of Formal Representation}

In the liminal space between phenomenological experience and formal representation lies a profound epistemological challenge: how do we capture the invariant structures of reality while acknowledging the perspectival nature of conceptualization? The edifice of computational ontology emerges from this tension, offering a systematic approach to formalizing our understanding of domains while maintaining critical awareness of the gap between representation and reality.

The genealogy we trace here—from extensional structures through conceptualizations to ontological commitments—mirrors a fundamental dialectic in systems theory itself. Von Bertalanffy's seminal characterization of a system as ``a complex of interacting elements'' \cite{Bertalanffy1968} gestures toward this same problematic: the relationship between the elements we identify (extensional) and the interactions that constitute systemhood (intensional). As we shall see, this distinction ramifies throughout the formal apparatus of ontology construction.